\documentclass[11pt]{article}
\usepackage{amsmath,amssymb,amsthm,mathtools,graphicx}
\usepackage[margin=1in]{geometry}
\usepackage[colorlinks=true,linkcolor=blue,citecolor=blue,urlcolor=blue]{hyperref}

\newtheorem{theorem}{Theorem}
\newtheorem{lemma}[theorem]{Lemma}
\newtheorem{proposition}[theorem]{Proposition}
\newtheorem{corollary}[theorem]{Corollary}
\theoremstyle{remark}
\newtheorem{remark}{Remark}

\newcommand{\Q}{\mathbb Q}
\newcommand{\R}{\mathbb R}
\newcommand{\e}{\mathrm e}
\newcommand{\diag}{\operatorname{diag}}
\newcommand{\Th}{\Theta}

\title{Solution to Ramanujan Challenge Problem 2.5:\\
Efficient Rational Approximation of Catalan's Constant~$G$}
\author{Xiang Huang}
\date{August 2026}

\begin{document}
\maketitle

\begin{abstract}
We prove, for every one of the three terminal columns, that the
matrix quotient in Problem~2.5 tends to Catalan's constant.  The
paper contains two exact proofs.  The first relates a Meijer--$G$
trajectory to the printed matrix by an explicit inverse-transpose
identity.  The second is independent of that special-function
connection: an elementary positive triple integral produces a fast
adjoint solution, while a rational companion and a contracting Miller
cone select its projective limit.  We also give the exact inverse
Delannoy transform and the requested finite integral formula for its
coefficient error.  No finite-prefix guess or decimal recognition is
used.
\end{abstract}

\section{The matrix and the assertion}

Let $M(n)=(m_{ij}(n))_{1\leq i,j\leq3}$ be the matrix in the
question.  For reference, its entries are
\begingroup\small
\begin{align*}
m_{11}&=(-2n-5)(n+3)^2(136n^4+1424n^3+5548n^2+9551n+6141),\\
m_{12}&=384n^6+6384n^5+44168n^4+162698n^3+336377n^2+369933n+169011,\\
m_{13}&=-480n^4-4980n^3-19210n^2-32690n-20730,\\
m_{21}&=(n+2)^2(n+3)^2(4n+10)(48n^3+386n^2+1017n+879),\\
m_{22}&=(n+2)^2(-272n^5-3848n^4-21732n^3-61184n^2-85761n-47808),\\
m_{23}&=(n+2)^2(320n^3+2540n^2+6610n+5640),\\
m_{31}&=(-4n-10)(n+2)^2(n+3)^2(32n^4+302n^3+1037n^2+1530n+813),\\
m_{32}&=(n+2)^2(192n^6+2984n^5+19116n^4+64452n^3+120256n^2+117279n+46476),\\
m_{33}&=(n+2)^2(-16n^5-408n^4-2912n^3-8884n^2-12254n-6240).
\end{align*}
\endgroup
Put
\[
 \mathcal M_N=M(0)M(1)\cdots M(N-1),\qquad
 A=\begin{pmatrix}30921&-32972&8240\\33750&-36000&9000\end{pmatrix},
\]
where the empty product is the identity, and define
\[
 A\mathcal M_N=
 \begin{pmatrix}P_{N,1}&P_{N,2}&P_{N,3}\\
                 Q_{N,1}&Q_{N,2}&Q_{N,3}\end{pmatrix}.
\]
Finally,
\[
 G=\sum_{r=0}^{\infty}\frac{(-1)^r}{(2r+1)^2}.
\]

\begin{theorem}[Main assertion]\label{thm:main-statement}
For $j=1,2,3$,
\[
       \lim_{N\to\infty}\frac{P_{N,j}}{Q_{N,j}}=G.
\]
\end{theorem}

The proof is completed in Theorem~\ref{thm:all-columns}.

\begin{lemma}[Catalan moment]\label{lem:catalan-moment}
The improper integral is absolutely convergent and
\[
 G=\int_0^1\frac{-\log t}{1+t^2}\,dt,
 \qquad \int_0^1(-\log t)\,dt=1.
\]
\end{lemma}

\begin{proof}
The second identity follows by integration by parts.  For $0\leq t<1$,
\[
 \frac1{1+t^2}=\sum_{r\geq0}(-1)^rt^{2r}.
\]
The partial sums have absolute value at most $2$, while
$-\log t$ is integrable on $(0,1)$.  Dominated convergence and
$\int_0^1t^{2r}(-\log t)\,dt=(2r+1)^{-2}$ give the first identity.
\end{proof}

The following exact moment identity will also be used in the
Delannoy appendix.  It is not used to infer decay merely from a
finite alternating sum.

\begin{lemma}[Exact moment lift]\label{lem:moment-lift}
Let $\mathbf p_0$ and $\mathbf q_0$ be the two rows of $A$, set
\[
 \mathbf R_0(X)=\mathbf q_0-(1+X)\mathbf p_0,
 \qquad R_{N,j}(X)=(\mathbf R_0(X)\mathcal M_N)_j.
\]
Then $R_{N,j}(X)=Q_{N,j}-(1+X)P_{N,j}$ and
\begin{align}
 Q_{N,j}&=\int_0^1(10-18x)R_{N,j}(x)\,dx,\label{eq:q-moment}\\
 P_{N,j}&=\int_0^1(6-12x)R_{N,j}(x)\,dx,\label{eq:p-moment}\\
 GQ_{N,j}-P_{N,j}
 &=\int_0^1\frac{-\log t}{1+t^2}R_{N,j}(t^2)\,dt.\label{eq:error-moment}
\end{align}
\end{lemma}

\begin{proof}
The displayed affine formula for $R_{N,j}$ follows directly from
the definition.  If $R(X)=a+bX$, then the four elementary moments
\[
 \int_0^1(10-18x)\,dx=1,\quad
 \int_0^1x(10-18x)\,dx=-1,\quad
 \int_0^1(6-12x)\,dx=0,\quad
 \int_0^1x(6-12x)\,dx=-1
\]
give \eqref{eq:q-moment} and \eqref{eq:p-moment}.  Lemma
\ref{lem:catalan-moment} gives
\[
 \int_0^1\frac{-\log t}{1+t^2}
       \bigl(Q-(1+t^2)P\bigr)\,dt=GQ-P,
\]
which is \eqref{eq:error-moment}.  The logarithmic endpoint is
integrable, so these are ordinary improper Lebesgue integrals.
\end{proof}

\section{Inverse-transpose selection}

We record the standard form of the recessive-solution principle
needed below.  It is a consequence of the Poincar\'e--Perron theorem
in its higher-order (discrete Levinson) form; see
\cite{WeinbaumCMF}.

\begin{proposition}[Recessive recognition and Miller selection]
\label{prop:miller}
Suppose a row companion system $X_{n+1}=X_nC(n)$ satisfies
$C(n)=C_\infty+O(n^{-1})$, where $C_\infty$ is invertible and
diagonalizable.  Suppose $C_\infty$ has a simple eigenvalue $\mu$
such that
\[
 0<|\mu|<\min\{|\lambda|:\lambda\in\operatorname{spec}(C_\infty),
 \ \lambda\ne\mu\}.
\]
If a nonzero scalar solution $F_n$ has
$\limsup |F_n|^{1/n}\leq|\mu|$, it spans the unique recessive
solution line.  Writing
\[
 \Phi_C(N)=C(0)\cdots C(N-1),\qquad
 \mathbf F_0=(F_0,\ldots,F_{r-1})^T,
\]
one has
\[
 [\Phi_C(N)^{-T}w_N]\longrightarrow[\mathbf F_0]
\]
for every terminal vector $w_N$ whose coefficient in the
corresponding dominant dual direction is nonzero.

Moreover, if $T(n)U(n+1)=U(n)C(n)$ with rational invertible $U(n)$,
then
\begin{equation}\label{eq:coboundary-selection}
 [\Phi_T(N)^{-T}e_j]\longrightarrow[U(0)^{-T}\mathbf F_0]
\end{equation}
whenever $U(N)^Te_j$ has nonzero recessive-dual connection.
\end{proposition}

\begin{proof}
The spectral gap splits the asymptotic solution space into the
one-dimensional $\mu$-space and its complementary block.  Discrete
Levinson asymptotics show that every solution outside the first
space has strictly larger exponential rate.  This proves the
recognition statement.  Inverting and transposing a fundamental
product replaces each exponential factor $\lambda^N$ by
$\lambda^{-N}$; hence the recessive primal line becomes the unique
dominant dual line.  This proves the projective assertion.

Finally, multiplication of the coboundary identities gives
\[
 \Phi_T(N)=U(0)\Phi_C(N)U(N)^{-1},\qquad
 \Phi_T(N)^{-T}=U(0)^{-T}\Phi_C(N)^{-T}U(N)^T.
\]
Rational gauges have at most polynomial growth along the positive
integers and therefore do not alter an exponentially isolated line.
This gives \eqref{eq:coboundary-selection}.
\end{proof}

\section{An exact Meijer--\texorpdfstring{$G$}{G} trajectory}

We use the Barnes convention of \cite[Sec.~16.17]{DLMF},
\[
G_{p,q}^{m,r}\!\left(z\,\middle|\,
 \begin{matrix}a_1,\ldots,a_p\\b_1,\ldots,b_q\end{matrix}\right)
=\frac1{2\pi i}\int_L
 \frac{\prod_{j=1}^{m}\Gamma(b_j-s)
       \prod_{j=1}^{r}\Gamma(1-a_j+s)}
      {\prod_{j=m+1}^{q}\Gamma(1-b_j+s)
       \prod_{j=r+1}^{p}\Gamma(a_j-s)}z^s\,ds.
\]
For $n\geq0$ define
\begin{equation}\label{eq:Fn-def}
 F_n(z)=G_{4,4}^{4,2}\!\left(z\,\middle|\,
 \begin{matrix}-1-n,-1-n,n+\frac52,n+3\\
 \frac32,\frac12,1,0\end{matrix}\right),
 \qquad F_n=F_n(1),
\end{equation}
and set $\Th=z\,d/dz$ and
\[
 Y_n=\bigl(F_n,(\Th F_n)(1),(\Th^2F_n)(1)\bigr).
\]

\begin{lemma}[Differential reduction]\label{lem:theta-reduction}
At $z=1$,
\begin{equation}\label{eq:theta3}
 (\Th^3F_n)(1)=
 -\frac{(n+2)^3(2n+3)}7F_n
 +\frac{2n^2+8n+5}{14}(\Th F_n)(1)
 +\frac{8n^2+30n+39}{14}(\Th^2F_n)(1).
\end{equation}
\end{lemma}

\begin{proof}
The Meijer--$G$ differential equation for \eqref{eq:Fn-def} is
\[
 \left[z(\Th+n+2)^2(\Th-n-\tfrac32)(\Th-n-2)
 -(\Th-\tfrac32)(\Th-\tfrac12)(\Th-1)\Th\right]F_n(z)=0.
\]
At $z=1$ the fourth-order terms cancel.  Expansion and division by
$7/2$ gives \eqref{eq:theta3}.
\end{proof}

Define $T(n)=(t_{ij}(n))$ by
\begin{center}
\resizebox{\textwidth}{!}{$
\begin{pmatrix}
\dfrac{4(n+2)(17n^3+111n^2+240n+171)}{(n+1)(n+3)(2n+3)(2n+5)}&
\dfrac{(n+2)(24n^2+101n+102)}{(n+1)(2n+3)}&
\dfrac{(n+2)(2n+5)(16n^2+81n+90)}{2(n+1)(2n+3)}\\[6pt]
\dfrac{96n^4+780n^3+2384n^2+3273n+1723}{(n+1)(n+2)(n+3)(2n+3)(2n+5)}&
\dfrac{68n^3+398n^2+778n+523}{2(n+1)(n+2)(2n+3)}&
\dfrac{96n^4+884n^3+2970n^2+4360n+2403}{4(n+1)(n+2)(2n+3)}\\[6pt]
\dfrac{-5(24n^2+117n+143)}{(n+1)(n+2)(n+3)(2n+3)(2n+5)}&
\dfrac{-5(16n+41)}{2(n+1)(n+2)(2n+3)}&
\dfrac{8n^3-44n^2-478n-801}{4(n+1)(n+2)(2n+3)}
\end{pmatrix}.
$}
\end{center}

\begin{lemma}[Exact contiguous trajectory]\label{lem:trajectory}
For every integer $n\geq0$,
\[
                         Y_{n+1}=Y_nT(n).
\]
\end{lemma}

\begin{proof}
In the Barnes integrand for $F_n$, the $n$-dependent factor is
\[
 \frac{\Gamma(n+2+s)^2}
      {\Gamma(n+\frac52-s)\Gamma(n+3-s)}.
\]
Its quotient after $n\mapsto n+1$ is
\[
 \frac{(n+2+s)^2}{(n+\frac52-s)(n+3-s)}.
\]
Apply this identity also after multiplication by $s$ and $s^2$.
Gamma recurrence and the cubic reduction \eqref{eq:theta3} reduce
the three resulting Barnes integrals to the basis $1,s,s^2$; the
three coefficient columns are exactly the displayed columns of
$T(n)$.  All operations are identities in $\Q(n,s)$ after the
gamma factors are cancelled.  They are first valid on a Barnes
contour for $0<z<1$; gamma decay permits contour shifting and
differentiation.  Dominated passage to $z=1$ gives the assertion.
Thus this is an exact contiguous calculation, not a check of
finitely many values.
\end{proof}

Put
\begin{equation}\label{eq:sigma}
 \sigma(n)=-2(n+2)^2(n+3)^2(2n+5)(2n+7)^2.
\end{equation}

\begin{lemma}[Inverse-transpose certificate]\label{lem:inverse-transpose}
For the printed matrix $M(n)$ and the displayed $T(n)$,
\begin{equation}\label{eq:single-certificate}
              T(n)^TM(n)=\sigma(n)I_3.
\end{equation}
Consequently, if $\Phi_T(N)=T(0)\cdots T(N-1)$ and
$H_N=\prod_{m=0}^{N-1}\sigma(m)$, then
\begin{equation}\label{eq:product-certificate}
 \mathcal M_N=H_N\Phi_T(N)^{-T},\qquad
 H_N=(-16)^N(2)_N^2(3)_N^2(\tfrac52)_N(\tfrac72)_N^2.
\end{equation}
\end{lemma}

\begin{proof}
Substitution of the nine displayed rational functions and the nine
polynomials defining $M(n)$ gives \eqref{eq:single-certificate}
after clearing the nonzero denominator
$(n+1)(n+2)(n+3)(2n+3)(2n+5)$.  This is a nine-entry polynomial
identity in $\Q[n]$.  Multiplying it for $0\leq n<N$ gives the
first identity in \eqref{eq:product-certificate}; factoring each
linear term in \eqref{eq:sigma} gives the stated Pochhammer product.
\end{proof}

\section{Companion form and exact connection data}

Let $e_1=(1,0,0)^T$ and define the Krylov matrix
\begin{equation}\label{eq:U-def}
 U(n)=\bigl[e_1,\ T(n)e_1,\ T(n)T(n+1)e_1\bigr].
\end{equation}

\begin{lemma}[Krylov determinant]\label{lem:U-det}
Let
\[
 D(n)=3072n^7+61824n^6+531184n^5+2525522n^4+7175793n^3
 +12183785n^2+11446123n+4589916.
\]
Then
\[
 \det U(n)=
 -\frac{15D(n)}{2(n+1)(n+2)^3(n+3)^2(n+4)
 (2n+3)^2(2n+5)^2(2n+7)}.
\]
In particular, $U(n)$ is invertible for every integer $n\geq0$.
\end{lemma}

\begin{proof}
Expand the determinant using \eqref{eq:U-def} and factor.  Every
coefficient of $D$ is positive and all denominator factors are
positive for $n\geq0$.
\end{proof}

\begin{lemma}[Companion coboundary]\label{lem:companion}
The matrix
\[
 C(n)=U(n)^{-1}T(n)U(n+1)
\]
is a row companion matrix, $T(n)U(n+1)=U(n)C(n)$, and
\[
 C(n)=C_\infty+O(n^{-1}),\qquad
 C_\infty=\begin{pmatrix}0&0&1\\1&0&-35\\0&1&35\end{pmatrix}.
\]
Moreover
\begin{equation}\label{eq:YU}
                 Y_nU(n)=(F_n,F_{n+1},F_{n+2}).
\end{equation}
\end{lemma}

\begin{proof}
The first two columns of the coboundary identity follow from the
definition of $U(n)$; the last column defines the scalar recurrence,
so $C(n)$ is companion.  Exact rational simplification gives the
displayed limit.  Equation \eqref{eq:YU} follows from
Lemma~\ref{lem:trajectory}.
\end{proof}

\begin{lemma}[Spectral gap]\label{lem:spectral-gap}
The characteristic polynomial of $C_\infty$ is
\[
 (\lambda-1)(\lambda^2-34\lambda+1).
\]
Its roots are
\[
 \mu=17-12\sqrt2=(\sqrt2-1)^4,\qquad 1,\qquad
 \mu^{-1}=17+12\sqrt2.
\]
Thus $\mu$ is simple and has strictly smallest modulus.  A right
eigenvector for $\mu$ is
\begin{equation}\label{eq:rmu}
 r_\mu=(17+12\sqrt2,-18-12\sqrt2,1)^T.
\end{equation}
\end{lemma}

\begin{proof}
This is direct factorization and multiplication by $C_\infty$.
\end{proof}

\begin{lemma}[All three terminal connections are nonzero]
\label{lem:terminal-connections}
For each $j=1,2,3$, the vector $U(N)^Te_j$ has nonzero connection
with the recessive dual mode in Proposition~\ref{prop:miller}.
\end{lemma}

\begin{proof}
Exact leading terms of the displayed rational matrix $U(n)$ give
\begin{equation}\label{eq:U-limit}
 \lim_{n\to\infty}\diag(1,n,n^3)U(n)
 =L:=\begin{pmatrix}1&17&577\\0&24&816\\0&-30&-1020\end{pmatrix}.
\end{equation}
The three row pairings of $L$ with \eqref{eq:rmu} are
\[
 96(3-2\sqrt2),\qquad 96(4-3\sqrt2),\qquad
 120(-4+3\sqrt2).
\]
None is zero: otherwise $\sqrt2$ would be rational.  Polynomial row
scalings do not change whether the corresponding asymptotic
coefficient vanishes.  Hence all three terminal vectors satisfy the
noncancellation hypothesis.
\end{proof}

Let $\mathbf F=(F_0,F_1,F_2)^T$.

\begin{lemma}[Barnes initial values]\label{lem:initial-F}
The initial vector is
\begin{equation}\label{eq:F-vector}
\mathbf F=\sqrt\pi\begin{pmatrix}
150G-128\log2-\dfrac{146}{3}\\[3pt]
\dfrac{24745}{4}G-\dfrac{14624}{3}\log2-\dfrac{823511}{360}\\[3pt]
\dfrac{7225281}{32}G-\dfrac{886784}{5}\log2
 -\dfrac{2818419551}{33600}
\end{pmatrix}.
\end{equation}
\end{lemma}

\begin{proof}
Insert $n=0,1,2$ in the Barnes integrand in
\eqref{eq:Fn-def}.  Cancelling the two gamma quotients gives
\[
 \frac{\Gamma(\frac12-s)\Gamma(-s)\Gamma(n+2+s)^2}
 {(\frac32-s)_{n+1}(1-s)_{n+2}}.
\]
Residue expansion (first at $z=r$ with $0<r<1$ and then with
$r\uparrow1$) reduces each of the three cases by partial
fractions to
\[
 \sum_{k\geq0}\frac{(-1)^k}{2k+1}=\frac\pi4,\quad
 \sum_{k\geq1}\frac{(-1)^{k-1}}k=\log2,\quad
 \sum_{k\geq0}\frac{(-1)^k}{(2k+1)^2}=G,
\]
plus finite rational tails.  Collecting the coefficients gives
\eqref{eq:F-vector}.  This is the exact Barnes-moment evaluation
the regulator
$r$ also justifies the contour endpoint before dominated passage to
$r=1$.
\end{proof}

\begin{lemma}[Initial coboundary]\label{lem:initial-U}
One has
\[
 U(0)=\begin{pmatrix}
 1&\frac{152}{5}&\frac{195477}{175}\\
 0&\frac{1723}{90}&\frac{1963751}{2800}\\
 0&-\frac{143}{18}&-\frac{165201}{560}
 \end{pmatrix}
\]
and
\begin{equation}\label{eq:Atilde}
 AU(0)^{-T}=\frac1{382493}
 \begin{pmatrix}
 30102216645&-19896711516&525137760\\
 32863650750&-21712357800&573048000
 \end{pmatrix}.
\end{equation}
\end{lemma}

\begin{proof}
Substitute $n=0$ in \eqref{eq:U-def}; exact inversion and
multiplication by $A$ give \eqref{eq:Atilde}.
\end{proof}

\begin{lemma}[Exact Catalan cancellation]\label{lem:cancellation}
The logarithmic and rational terms cancel exactly:
\begin{equation}\label{eq:exact-cancellation}
            AU(0)^{-T}\mathbf F=450\sqrt\pi\binom{G}{1}.
\end{equation}
\end{lemma}

\begin{proof}
Multiply the rational matrix \eqref{eq:Atilde} by
\eqref{eq:F-vector}.  The two coefficients of $\log2$ are zero,
the two rational constant terms are zero, and the remaining
coefficients are respectively $450G$ and $450$.  This is an identity
in the three-dimensional $\Q$-space spanned by $1,\log2,G$; no
decimal recognition is involved.
\end{proof}

\section{Exact Delannoy transform}

This section repairs the finite-prefix claims in the earlier
Delannoy argument.  It is included because the transform and its
error integral are useful exact identities.  The main convergence
proof below does not require a sign assertion about the transformed
coefficient $f(k)$.

Let
\[
 \mathcal B(N,k)=2^k\binom{2k}{k}\binom Nk\binom{N+k}{k}
 \quad(0\leq k\leq N).
\]

\begin{lemma}[Inverse Delannoy matrix]\label{lem:inverse-delannoy}
The inverse of the lower-triangular matrix $\mathcal B$ is
\begin{equation}\label{eq:inverse-delannoy}
 \beta(k,N)=
 \frac{(-1)^{k+N}(2N+1)\binom{k}{N}}
 {2^k\binom{2k}{k}(k+N+1)\binom{k+N}{N}}
 \qquad(0\leq N\leq k).
\end{equation}
\end{lemma}

\begin{proof}
We prove
$\sum_{N=j}^k\beta(k,N)\mathcal B(N,j)=\delta_{kj}$.
For $k=j$ the sole product is $1$.  Suppose $m=k-j>0$ and put
$N=j+r$.  After removing a nonzero factor independent of $r$ and
using
\[
 \frac{2j+2r+1}{k+j+r+1}
 =2-\frac{2k+1}{k+j+r+1},
\]
the sum becomes
\begin{align*}
&2\,{}_2F_1(-m,2j+1;k+j+1;1)\\
&\quad-\frac{2k+1}{k+j+1}
 {}_2F_1(-m,2j+1;k+j+2;1).
\end{align*}
Chu--Vandermonde evaluates the two hypergeometric terms as
\[
 \frac{(m)_m}{(k+j+1)_m},\qquad
 \frac{(m+1)_m}{(k+j+2)_m}.
\]
Their ratio is $2(k+j+1)/(2k+1)$, so the displayed difference is
zero.  This proves the inverse formula for every index, rather than
checking a finite square.
\end{proof}

Define the scalar normalization
\[
 H_N=(-16)^N(2)_N^2(3)_N^2(\tfrac52)_N(\tfrac72)_N^2
\]
and, for the first column only in this section,
\[
 \widehat P_N=P_{N,1}/H_N,\qquad
 \widehat Q_N=Q_{N,1}/H_N.
\]
There are unique rational sequences $f,g$ such that
\begin{equation}\label{eq:delannoy-expansion}
 \widehat Q_N=\sum_{k=0}^Nf(k)\mathcal B(N,k),\qquad
 \widehat P_N=\sum_{k=0}^Ng(k)\mathcal B(N,k),
\end{equation}
namely
\begin{equation}\label{eq:inverse-sums}
 f(k)=\sum_{N=0}^k\beta(k,N)\widehat Q_N,\qquad
 g(k)=\sum_{N=0}^k\beta(k,N)\widehat P_N.
\end{equation}

\begin{theorem}[Holonomic $k$-recurrences]\label{thm:k-closure}
The sequences $f(k)$ and $g(k)$ in \eqref{eq:inverse-sums} are
P-recursive.  This conclusion holds for all $k$ as a consequence of
holonomic closure applied to the exact finite sums; it is not based
on a guessed operator or verification of a finite prefix.
\end{theorem}

\begin{proof}
By \eqref{eq:product-certificate}, $\widehat P_N$ and
$\widehat Q_N$ are coordinates of a fixed row multiplied by the
three-dimensional rational first-order system $\Phi_T(N)^{-T}$.
Four consecutive coordinate functionals are linearly dependent over
$\Q(N)$; clearing denominators gives a polynomial-coefficient scalar
recurrence.  Thus both sequences are holonomic in $N$.

The kernel $\beta(k,N)$ is proper hypergeometric: both of its shift
quotients are rational functions of $(k,N)$.  Zeilberger's creative
telescoping theorem \cite{Zeilberger}, equivalently closure of holonomic functions
under definite summation, applied to the finite range $0\leq N\leq k$
therefore produces polynomial-coefficient recurrences for the two
sums \eqref{eq:inverse-sums}.  The argument proves existence and
annihilation identically; no order, degree, or minimality inferred by
guessing is asserted.
\end{proof}

\begin{theorem}[Effective positivity and recessiveness]\label{thm:positive-recessive}
For every $n\geq0$, $F_n>0$, and
\begin{equation}\label{eq:Fn-rate}
             \lim_{n\to\infty}F_n^{1/n}=17-12\sqrt2=\mu.
\end{equation}
Consequently $(F_n)$ is the unique recessive solution of the
companion recurrence in Lemma~\ref{lem:companion}.
\end{theorem}

\begin{proof}
Two beta integrals, two Laplace integrals, and the duplication formula
transform the Barnes integral exactly into
\begin{align}
F_n=\frac{\sqrt\pi}{\Gamma(n+1)\Gamma(n+2)}
\int_{(0,1)^2\times(0,\infty)^2}
 &x^{1/2}(1-x)^n(1-y)^{n+1}u^{n+1}v^{n+1}\notag\\[-2pt]
 &\times\exp\!\left(-u-v-2\sqrt{\frac{uv}{xy}}\right)
 \,du\,dv\,dx\,dy.\label{eq:positive-F}
\end{align}
Every factor is positive, proving $F_n>0$ directly for every $n$.

For the exact root rate put $u=na$, $v=nb$ and write
\[
 F_n=\frac{\sqrt\pi\,n^{2n+4}}{\Gamma(n+1)\Gamma(n+2)}J_n,\qquad
 J_n=\int_D A\e^{n\Phi}\,dx\,dy\,da\,db,
\]
where $D=(0,1)^2\times(0,\infty)^2$,
\[
 A=x^{1/2}(1-y)ab
\]
and
\[
 \Phi=\log(1-x)+\log(1-y)+\log a+\log b-a-b
       -2\sqrt{ab/(xy)}.
\]
Let $p=\sqrt{xy}$ and $r=\sqrt{ab}$.  Arithmetic--geometric mean
gives
\[
 \Phi\leq2\log(1-p)+2\log r-2r(1+p^{-1}),
\]
with equality only when $x=y$ and $a=b$.  Maximizing first over $r$
and then over $p$ yields
\[
 \Phi+2\leq2\log\frac{p(1-p)}{1+p}
 \leq4\log(\sqrt2-1)=:\Phi_*+2.
\]
The unique maximizer is
\[
 x=y=\sqrt2-1,\qquad a=b=\frac{\sqrt2-1}{\sqrt2}.
\]
Moreover
\[
 A\e^\Phi\leq
 x^{1/2}(1-x)(1-y)^2a^2b^2\e^{-a-b},
\]
and the right-hand side is integrable on $D$.  Hence
$J_n\leq\e^{(n-1)\Phi_*}\int_DA\e^\Phi$.  Conversely, on a fixed
small neighborhood of the unique maximizer, $A$ is bounded below by
a positive constant and $\Phi\geq\Phi_*-\varepsilon$.  It follows
that $J_n^{1/n}\to\e^{\Phi_*}$; this is the elementary Laplace
principle with explicit maximizer, not a finite numerical check.
Stirling's formula gives
\[
 \left(\frac{\sqrt\pi\,n^{2n+4}}
 {\Gamma(n+1)\Gamma(n+2)}\right)^{1/n}\longrightarrow\e^2.
\]
Therefore
\[
 \lim F_n^{1/n}=\e^{\Phi_*+2}=(\sqrt2-1)^4=17-12\sqrt2.
\]
Lemma~\ref{lem:spectral-gap} and Proposition~\ref{prop:miller}
now identify $(F_n)$ as the unique recessive solution.
\end{proof}

\begin{theorem}[Exact integral for the Delannoy coefficient error]
\label{thm:coefficient-error}
Put $h(k)=g(k)-Gf(k)$ and
\begin{equation}\label{eq:Psi-def}
 \Psi_k(t)=\sum_{N=0}^k\beta(k,N)\frac{R_{N,1}(t^2)}{H_N}.
\end{equation}
Then, for every $k\geq0$,
\begin{equation}\label{eq:h-integral-final}
 h(k)=-\int_0^1\frac{-\log t}{1+t^2}\Psi_k(t)\,dt.
\end{equation}
The sum in \eqref{eq:Psi-def} is finite and has rational
coefficients; in fact
\begin{equation}\label{eq:Psi-affine}
                 \Psi_k(t)=f(k)-(1+t^2)g(k).
\end{equation}
\end{theorem}

\begin{proof}
Divide \eqref{eq:error-moment} by $H_N$, multiply by
$\beta(k,N)$, and sum over $0\leq N\leq k$.  The sum has only
$k+1$ terms, so its interchange with the absolutely convergent
integral is automatic.  Equations \eqref{eq:inverse-sums} give
\eqref{eq:h-integral-final}.  Applying the same inverse sum to
$R_{N,1}(t^2)/H_N=\widehat Q_N-(1+t^2)\widehat P_N$ gives
\eqref{eq:Psi-affine}.
\end{proof}

\begin{remark}
Equation \eqref{eq:Psi-affine} shows why the integral identity alone
must not be advertised as a decay estimate: without an independent
bound it is an exact rewriting of $h(k)$.  The earlier inference from
finite cancellation to an all-$k$ $O(8^{-k})$ bound is therefore
deleted.  The positive integral \eqref{eq:positive-F} supplies the
effective all-index argument actually needed for the proof.
\end{remark}

\begin{remark}[Other audit corrections]
If the unused raw scalar recurrence is indexed as
\[
 \sum_{j=0}^3c_j(n)y_{n+j}=0,
\]
the degree pattern is
\[
 (28,21,14,7),\qquad \deg c_j=7(4-j),
\]
not $\deg c_j=7(3-j)$.
The former modular paragraph has also been removed: the eta product
$\eta^2(2\tau)\eta^2(4\tau)$ has formal weight $2$, and no modular
claim is required here.  All endpoint passages used above are made
by dominated convergence or by the positive Laplace integral.
\end{remark}

\section{Limit identification and all three columns}

\begin{theorem}[Hatted sequences and the three-column bridge]
\label{thm:all-columns}
For every $N\geq0$ and $j=1,2,3$, define
\begin{equation}\label{eq:hats-all}
 \widehat P_{N,j}=\frac{P_{N,j}}{H_N},\qquad
 \widehat Q_{N,j}=\frac{Q_{N,j}}{H_N},
 \qquad
 H_N=(-16)^N(2)_N^2(3)_N^2(\tfrac52)_N(\tfrac72)_N^2.
\end{equation}
Then
\[
 \frac{P_{N,j}}{Q_{N,j}}
 =\frac{\widehat P_{N,j}}{\widehat Q_{N,j}}\longrightarrow G
 \qquad(j=1,2,3).
\]
\end{theorem}

\begin{proof}
By \eqref{eq:product-certificate}, for each standard column $e_j$,
\begin{equation}\label{eq:hat-vector}
 \binom{\widehat P_{N,j}}{\widehat Q_{N,j}}
                   =A\Phi_T(N)^{-T}e_j.
\end{equation}
Theorem~\ref{thm:positive-recessive} identifies $\mathbf F$ as the
recessive initial vector.  Lemma~\ref{lem:terminal-connections}
verifies the nonzero terminal coefficient separately for each of
$j=1,2,3$.  Proposition~\ref{prop:miller} and
Lemma~\ref{lem:cancellation} therefore give
\[
 \left[A\Phi_T(N)^{-T}e_j\right]
 \longrightarrow
 \left[AU(0)^{-T}\mathbf F\right]
 =\left[450\sqrt\pi\binom G1\right]=\left[\binom G1\right].
\]
The second coordinate of a normalized representative tends to the
nonzero number $1$, so $\widehat Q_{N,j}$ (and hence $Q_{N,j}$) is
nonzero for all sufficiently large $N$, and
$\widehat P_{N,j}/\widehat Q_{N,j}\to G$.  Finally $H_N\ne0$, so
the equality of the hatted and original ratios in
\eqref{eq:hats-all} proves all three assertions of
Theorem~\ref{thm:main-statement}.
\end{proof}

\section{An independent proof by an elementary Miller cone}
\label{sec:miller-independent}

This section gives a second proof of the main assertion.  It uses only
the printed matrix, elementary real integrals, rational identities, and a
uniform contraction argument.  In particular, it does not use a
Meijer--$G$ function, a Barnes integral, or a special-function connection
formula.

\subsection{The common projective limit}

Put $S=\operatorname{diag}(1,-1,1)$ and
\[
        B_n=-SM(n)S.
\]
Every entry of $B_n$ is positive for $n\geq0$.  If
$\boldsymbol p_n,\boldsymbol q_n$ are the original numerator and
denominator rows, define
\[
 \widetilde{\boldsymbol p}_n=(-1)^n\boldsymbol p_nS,\qquad
 \widetilde{\boldsymbol q}_n=(-1)^n\boldsymbol q_nS.
\]
Then
\[
  \widetilde{\boldsymbol p}_{n+1}
       =\widetilde{\boldsymbol p}_nB_n,\qquad
  \widetilde{\boldsymbol q}_{n+1}
       =\widetilde{\boldsymbol q}_nB_n,
\]
and the signs cancel in the quotient:
\[
 \frac{\widetilde p_{n,j}}{\widetilde q_{n,j}}
       =\frac{P_{n,j}}{Q_{n,j}}.
\]

The following elementary contraction will also be used at the end.

\begin{lemma}[Common-limit lemma]\label{lem:miller-common-limit}
There is a real number $L$ such that
\[
       \frac{P_{n,j}}{Q_{n,j}}\longrightarrow L
       \qquad(j=1,2,3).
\]
More precisely, all three ratios lie in an interval whose length is
$O((2/3)^n)$.
\end{lemma}

\begin{proof}
Set
\[
 X_n=(n+1)\frac{\widetilde q_{n,2}}{\widetilde q_{n,1}},
 \qquad
 Y_n=(n+1)^3\frac{\widetilde q_{n,3}}{\widetilde q_{n,1}}.
\]
Direct substitution of the entries of $B_n$ gives the invariant rectangle
\begin{equation}\label{eq:miller-den-cone}
          1\leq X_n\leq\frac32,\qquad 0\leq Y_n\leq2.
\end{equation}
Here is a short exact way to check the induction.  After the common positive
factor $\widetilde q_{n,1}$ is removed, each new coordinate is affine in
$(X_n,Y_n)$.  Each of the four required inequalities is therefore checked at
the four corners
\[
       (1,0),\ (1,2),\ (3/2,0),\ (3/2,2).
\]
After clearing the positive denominators, the resulting expressions are
polynomials in $n$ with nonnegative coefficients.

Let $r_{n,i}=\widetilde p_{n,i}/\widetilde q_{n,i}$.  Positivity of $B_n$
writes each $r_{n+1,j}$ as a convex combination of the three $r_{n,i}$.
Using \eqref{eq:miller-den-cone}, the normalized weights of both
$r_{n,1}$ and $r_{n,2}$ are at least $1/6$.  Thus, if
$\ell_n\leq r_{n,i}\leq u_n$ for all $i$, then
\[
\frac{r_{n,1}+r_{n,2}}6+\frac23\ell_n
 \leq r_{n+1,j}\leq
\frac{r_{n,1}+r_{n,2}}6+\frac23u_n .
\]
Defining the next enclosing interval by these two endpoints gives
$u_{n+1}-\ell_{n+1}=(2/3)(u_n-\ell_n)$.  The nested lower and upper
endpoints consequently have one common limit, and they squeeze every
$r_{n,j}$.
\end{proof}

\subsection{A positive adjoint moment}

For $p,q,v\in(0,1)$ put
\[
       D=pq(1+v^2)+2v
\]
and define
\begin{align}
 I_n(A,B,C;k)
   =\iiint_{(0,1)^3}
   \frac{16p^{2n+6+A}q^{2n+5+B}
      (1-p^2)^n(1-q^2)^{n+1}v^{2n+3+C}}
        {D^{\,2n+4+k}}\,dp\,dq\,dv .        \label{eq:miller-moment}
\end{align}
All these integrals are finite and positive.  Define
\begin{equation}\label{eq:miller-w}
\begin{split}
 w_{n,1}&=I_n(0,0,0;0),\\
 w_{n,2}&=2(n+2)I_n(0,0,1;1),\\
 w_{n,3}&=-(n+2)I_n(0,0,1;1)
       +2(n+2)(2n+5)I_n(0,0,2;2).
\end{split}
\end{equation}
The last component is also positive.  One proof is to integrate in $v$
after applying integration by parts to a multiple of
$v^{2n+4}(1-v^2)/D^{2n+5}$; the boundary terms vanish and the remaining
integrand is positive.

\begin{lemma}[Adjoint recurrence]\label{lem:miller-adjoint}
For
\[
 \lambda_n=(n+1)(n+2)^2(n+3)^2(2n+7)^2
\]
the vector in \eqref{eq:miller-w} satisfies
\begin{equation}\label{eq:miller-adjoint}
                 B_nw_{n+1}=\lambda_nw_n .
\end{equation}
Moreover
\begin{equation}\label{eq:miller-fast-bound}
             0<w_{n,1}\leq \frac{512}{64^n}.
\end{equation}
\end{lemma}

\begin{proof}
The recurrence has a three-variable divergence certificate.  To state the
certificate compactly, for a polynomial $H=H(n,p,q,v)$ define
\begin{align*}
\mathcal L_p(H)
={}&D\{p(1-p^2)H_p+[(2n+7)-(4n+9)p^2]H\}\\
 &-(2n+7)p(1-p^2)q(1+v^2)H,\\
\mathcal L_q(H)
={}&D\{q(1-q^2)H_q+[(2n+6)-(4n+10)q^2]H\}\\
 &-(2n+7)q(1-q^2)p(1+v^2)H,\\
\mathcal L_v(H)
={}&D\{v(1-v^2)H_v+[(2n+4)-(2n+6)v^2]H\}\\
 &-(2n+7)v(1-v^2)(2pqv+2)H.
\end{align*}
For each of the three rows there are explicit polynomials
$H_{i,p},H_{i,q},H_{i,v}\in\mathbb Z[n,p,q,v]$ such that, after the base
integrand in \eqref{eq:miller-moment} is factored out,
\begin{equation}\label{eq:miller-certificate}
 4(2n+3)(n+2)
 \left(\sum_{j=1}^3(B_n)_{ij}W_{n+1,j}
                  -\lambda_nW_{n,i}\right)
 =
 \mathcal L_p(H_{i,p})+\mathcal L_q(H_{i,q})
                         +\mathcal L_v(H_{i,v}).
\end{equation}
Here $W_{n,i}$ denotes the rational factor which produces the $i$th
integrand in \eqref{eq:miller-w}.  The sparse integer coefficient lists for
the nine polynomials are included in the accompanying Lean source.  Its
kernel checks \eqref{eq:miller-certificate} as a polynomial identity, not by
sampling.  Integration of the right side is zero: the three antiderivatives
contain, respectively, $p(1-p^2)$, $q(1-q^2)$, and $v(1-v^2)$, so both
endpoint terms vanish.  This proves \eqref{eq:miller-adjoint}.

For \eqref{eq:miller-fast-bound}, rewrite the first integrand as
\[
\frac{512}{64^n}\,p(1-q^2)
 [4p^2(1-p^2)]^n[4q^2(1-q^2)]^n
 \left(\frac{pq}{D}\right)^5
 \left(\frac{2v}{D}\right)^{2n+3}
 \left(\frac D4\right)^4 .
\]
Every factor after $512/64^n$ belongs to $[0,1]$ on the unit cube.
The cube has volume one, giving the required bound.
\end{proof}

At the initial index the same moments select Catalan's constant exactly.
Writing
\[
 \widetilde{\boldsymbol p}_0=(30921,32972,8240),\qquad
 \widetilde{\boldsymbol q}_0=(33750,36000,9000),
\]
one has
\begin{equation}\label{eq:miller-initial-pair}
       \widetilde{\boldsymbol q}_0\cdot w_0=75,\qquad
       \widetilde{\boldsymbol p}_0\cdot w_0=75G .
\end{equation}
These are elementary integral evaluations.  Expanding
\eqref{eq:miller-w}, the two coefficient triples are
\[
 (33750,126000,180000),\qquad
 (30921,115408,164800).
\]
Fubini reduction first gives $75/16$ and $75G/16$ before the outer factor
$16$ in \eqref{eq:miller-moment}.  The remaining one-variable substitution
$t=2v/(1+v^2)$ reduces the second evaluation to
\[
              G=\int_0^1\frac{-\log t}{1+t^2}\,dt.
\]
All endpoint terms are rational functions and logarithms; hence
\eqref{eq:miller-initial-pair} is an exact evaluation, not numerical
recognition.

\subsection{A rational companion and its Miller graph}

Let
\[
\sigma_n=-2(n+2)^2(n+3)^2(2n+5)(2n+7)^2,\qquad
        T_n=\sigma_nM(n)^{-T}.
\]
Thus $T_n^TM(n)=\sigma_nI$.  With $e_1=(1,0,0)^T$, form the Krylov matrix
\[
       K_n=[\,e_1,\ T_ne_1,\ T_nT_{n+1}e_1\,].
\]
Direct rational simplification gives
\begin{equation}\label{eq:miller-companion}
       T_nK_{n+1}=K_nC_n,\qquad
 C_n=\begin{pmatrix}0&0&c_{0,n}\\1&0&c_{1,n}\\0&1&c_{2,n}\end{pmatrix}.
\end{equation}
For completeness, the coefficients are
\begin{align}
c_{0,n}&=\frac{(n+2)^2(2n+7)d(n+1)}
 {(n+4)(n+5)(2n+9)d(n)},\nonumber\\
c_{1,n}&=-\frac{\pi(n)}
 {(n+4)(n+5)(2n+5)(2n+9)d(n)},\label{eq:miller-coefficients}\\
c_{2,n}&=\frac{\omega(n)}
 {4(n+3)(n+5)(2n+7)(2n+9)d(n)},\nonumber
\end{align}
where
\begin{align*}
d(n)={}&3072n^7+61824n^6+531184n^5+2525522n^4\\
 &+7175793n^3+12183785n^2+11446123n+4589916,\\
\pi(n)={}&430080n^{11}+15954432n^{10}+267573056n^9\\
&\quad+2677696216n^8\\
&+17764567836n^7+82030514302n^6+269008975593n^5\\
&+626475469357n^4+1015303508704n^3+1090544343906n^2\\
&+698695833276n+202285840752,\\
\omega(n)={}&1720320n^{11}+62195712n^{10}+1017208064n^9\\
&\quad+9933863520n^8\\
&+64362354608n^7+290489793800n^6+931925873076n^5\\
&+2125092069586n^4+3375535586927n^3+3557018693358n^2\\
&+2237964273360n+636907494879.
\end{align*}
Clearing the positive denominators in \eqref{eq:miller-coefficients} gives
polynomials with nonnegative coefficients and proves, for every $n\geq0$,
\begin{equation}\label{eq:miller-coeff-bounds}
 1\leq c_{0,n}\leq\frac43,\quad
 -49\leq c_{1,n}\leq-35,\quad
 35\leq c_{2,n}\leq37.
\end{equation}
Set
\[
 A_n=-\frac{c_{1,n}}{c_{0,n}},\qquad
 B_n'=-\frac{c_{2,n}}{c_{0,n}},\qquad
 R_n=\frac1{c_{0,n}}.
\]
Then
\begin{equation}\label{eq:miller-pull-bounds}
 \frac{105}{4}\leq A_n\leq49,\qquad
 -37\leq B_n'\leq-\frac{105}{4},\qquad
 \frac34\leq R_n\leq1.
\end{equation}

Let
\[
 {\cal K}=[0,1/16]\times[0,1/256]
\]
and define the backward projective map
\begin{equation}\label{eq:miller-pull}
 {\cal P}_n(u,v)=
 \left(\frac1{A_n+B_n'u+R_nv},
       \frac{u}{A_n+B_n'u+R_nv}\right).
\end{equation}
Equations \eqref{eq:miller-pull-bounds} imply
\[
 23\leq A_n+B_n'u+R_nv\leq50
       \qquad((u,v)\in{\cal K}).
\]
Consequently ${\cal P}_n({\cal K})\subseteq{\cal K}$, and a direct
difference estimate gives
\begin{equation}\label{eq:miller-contraction}
 \|{\cal P}_n(x)-{\cal P}_n(y)\|_\infty
       \leq\frac18\|x-y\|_\infty .
\end{equation}
Thus there is a unique infinite pullback graph
$z_n={\cal P}_n(z_{n+1})$ with $z_n\in{\cal K}$; truncating at depth $h$
changes $z_n$ by at most $8^{-h}$ times the diameter of ${\cal K}$.

It remains to show that the positive moment is precisely the solution on
this graph.  Put $r_n=Sw_n$ and
\[
 s_n=-\frac{\sigma_n}{\lambda_n}
       =\frac{2(2n+5)}{n+1}\leq10,\qquad
 a_0=1,\quad a_{n+1}=a_ns_n,\quad f_n=a_nw_{n,1}.
\]
The adjoint identity and \eqref{eq:miller-companion} imply that
\[
             F_n=(f_n,f_{n+1},f_{n+2})
\]
obeys the companion evolution $F_{n+1}=F_nC_n$.  Moreover
\begin{equation}\label{eq:miller-f-fast}
             |f_n|\leq512(5/32)^n.
\end{equation}

For any companion solution $X=(X_0,X_1,X_2)$ define its defect from the
graph by
\[
 \delta_n(X)=(X_1-z_{n,1}X_0,\ X_2-z_{n,2}X_0).
\]
Using $z_n={\cal P}_n(z_{n+1})$ and \eqref{eq:miller-companion},
\begin{align*}
\delta_{n+1,1}&=-z_{n+1,1}\delta_{n,1}+\delta_{n,2},\\
\delta_{n+1,2}&=(c_{1,n}-z_{n+1,2})\delta_{n,1}
                         c_{2,n}\delta_{n,2}.
\end{align*}
The bounds above give
\begin{equation}\label{eq:miller-defect-growth}
             \|\delta_{n+1}(X)\|_\infty
                  \geq\frac12\|\delta_n(X)\|_\infty.
\end{equation}
On the other hand \eqref{eq:miller-f-fast} gives
\[
 2^n\|\delta_n(F_n)\|_\infty
       \leq544(5/16)^n\longrightarrow0.
\]
Iterating \eqref{eq:miller-defect-growth} backwards forces
$\delta_0(F_0)=0$:
\begin{equation}\label{eq:miller-selected}
              F_0=f_0(1,z_{0,1},z_{0,2}).
\end{equation}
This is the elementary Miller selection theorem needed here.

\subsection{Initial selection and the literal terminal column}

The initial Krylov matrix and two fixed covectors are
\[
K_0=\begin{pmatrix}
1&152/5&195477/175\\
0&1723/90&1963751/2800\\
0&-143/18&-165201/560
\end{pmatrix},
\]
\[
\ell_P=\begin{pmatrix}30102216645\\-19896711516\\525137760\end{pmatrix},
\qquad
\ell_Q=\begin{pmatrix}32863650750\\-21712357800\\573048000\end{pmatrix}.
\]
Exact multiplication gives
\begin{equation}\label{eq:miller-covectors}
 K_0\ell_P=382493\,\boldsymbol p_0^T,\qquad
 K_0\ell_Q=382493\,\boldsymbol q_0^T.
\end{equation}
Together with \eqref{eq:miller-initial-pair} and
\eqref{eq:miller-selected}, this yields
\begin{equation}\label{eq:miller-graph-value}
 \frac{\ell_P\cdot(1,z_{0,1},z_{0,2})}
      {\ell_Q\cdot(1,z_{0,1},z_{0,2})}=G.
\end{equation}
The denominator is positive throughout ${\cal K}$; this is immediate by
checking the two relevant endpoints of its affine expression.

We finally connect the abstract pullback to an actual column of the
challenge product.  Put $A_n=C_n^{-T}$.  The identities already proved give
the exact telescope
\begin{equation}\label{eq:miller-telescope}
 K_0^T\mathcal M_N
   =\Sigma_N(A_0A_1\cdots A_{N-1})K_N^T
\end{equation}
for a nonzero scalar $\Sigma_N$.  For $N=h+2$, direct rational
simplification gives
\begin{equation}\label{eq:miller-terminal}
 A_hA_{h+1}K_N^Te_3
       =\kappa_N(1,\tau_N,0)^T,\qquad \kappa_N>0,
\end{equation}
where
\[
\begin{aligned}
V_N={}&N^2(2N+1)(2N+3)
 (48N^4+306N^3+715N^2+729N+275),\\
R_N={}&6528N^8+54480N^7+187144N^6+344616N^5\\
 &+371338N^4+239883N^3+91396N^2+19200N+1800,
\qquad \tau_N=\frac{V_N}{R_N}.
\end{aligned}
\]
Clearly $\tau_N\geq0$, while $\tau_N\leq1/16$ follows after clearing the
positive denominator from
\begin{align*}
&R_N-16V_N\\
&\quad=3456N^8+28752N^7+99912N^6+191752N^5+226106N^4\\
&\qquad\quad+169691N^3+78196N^2+19200N+1800\geq0,
\end{align*}
where $R_N$ is the denominator displayed in $\tau_N$.  Hence
$(\tau_N,0)\in{\cal K}$.

Applying \eqref{eq:miller-telescope} to $e_3$ and cancelling its nonzero
scalar factors gives
\[
\frac{P_{h+2,3}}{Q_{h+2,3}}
=
\frac{\ell_P\cdot
 (1,{\cal P}_0{\cal P}_1\cdots{\cal P}_{h-1}(\tau_{h+2},0))}
{\ell_Q\cdot
 (1,{\cal P}_0{\cal P}_1\cdots{\cal P}_{h-1}(\tau_{h+2},0))}.
\]
By \eqref{eq:miller-contraction}, its chart tends to $z_0$; by
\eqref{eq:miller-graph-value}, the quotient tends to $G$.
Lemma~\ref{lem:miller-common-limit} then promotes the third-column limit to
all three columns.  This proves the main assertion independently.

\paragraph{Formal verification.}
The accompanying Lean development proves every analytic estimate,
polynomial certificate, matrix identity, and limit used in this section.
Its final theorem is
\[
       \texttt{problem25\_solved : Problem25Claim}.
\]
It contains no \texttt{sorry}.  Its axiom audit is
\[
 \{\texttt{propext},\ \texttt{Classical.choice},\ \texttt{Quot.sound}\}.
\]


\begin{remark}[Provenance of the two proofs]
The Meijer--$G$ route above was reconstructed after publication of the
official solutions and uses their special-function trajectory.  The proof in
Section~\ref{sec:miller-independent} was developed independently from the
printed challenge matrix.  It uses no Meijer--$G$ or Barnes connection data.
\end{remark}

\begin{thebibliography}{9}

\bibitem{WeinbaumCMF}
S.~Weinbaum et al.,
\emph{On Conservative Matrix Fields: Continuous Asymptotics and Arithmetic},
arXiv:2507.08138, 2025.

\bibitem{DLMF}
NIST Digital Library of Mathematical Functions,
\emph{Meijer $G$-function}, Section~16.17,
\url{https://dlmf.nist.gov/16.17}.

\bibitem{Zeilberger}
D.~Zeilberger,
\emph{The method of creative telescoping},
J. Symbolic Comput. \textbf{11} (1991), 195--204.

\end{thebibliography}

\end{document}
